\documentclass[../../main.tex]{subfiles}% 注意这里的文件路径不能用 ./main.tex ,否则用latexmk编译子文件会报错
\graphicspath{{\subfix{./image/}}} % 指定图片目录,后续可以直接使用图片文件名
% 注意这里的文件路径不能用 ../../image/ ,否则用latexmk编译子文件会报错

% 例如:
% \begin{figure}[H]
% \centering
% \includegraphics[scale=0.3]{图.png}
% \caption{}
% \label{figure:图}
% \end{figure}
% 注意:上述\label{}一定要放在\caption{}之后,否则引用图片序号会只会显示??.

\begin{document}

\section{Riesz-Fredholm理论}

\begin{examples}
设$\mathscr{X}=\mathbb{R}^n$，$T\in\mathscr{L}(\mathbb{R}^n)$，算子方程
\begin{align}
Tx=y \label{eq:op-eq}
\end{align}
有解的充要条件为$\langle z,y\rangle=0$对所有满足$T^*z=\theta$的$z\in\mathbb{R}^n$成立，其中$T^*$为$T$的转置。
该算子方程仅有两种情形：
\begin{enumerate}[(1)]
\item 对任意$y\in\mathbb{R}^n$，方程\eqref{eq:op-eq}存在唯一解；
\item 齐次方程$Tx=\theta$存在非零解，此时$Tx=\theta$与$T^*x=\theta$的非零解的极大线性无关组的基数相等。
\end{enumerate}
\end{examples}

\begin{examples}
设$K\in C([0,1]\times[0,1])$，考察积分方程
\begin{align}
x(t)=\int_{0}^{1}K(t,s)x(s)\mathrm{d}s+y(t) \label{eq:int-eq1}
\end{align}
及其共轭方程
\begin{align}
f(t)=\int_{0}^{1}K(s,t)f(s)\mathrm{d}s+g(t) \label{eq:int-eq2}
\end{align}
其中$x,y,f,g\in L^2[0,1]$，则有如下结论：
\begin{enumerate}[(1)]
\item 方程\eqref{eq:int-eq1}仅有两种情形：对任意$y\in L^2[0,1]$存在唯一解，或齐次方程存在非零解；
\item 方程\eqref{eq:int-eq2}与\eqref{eq:int-eq1}情形一致，对应齐次方程的线性无关解个数相同且有限；
\item 第二种情形下，方程\eqref{eq:int-eq1}有解当且仅当$\int_{0}^{1}f(t)y(t)\mathrm{d}t=0$（$f$为\eqref{eq:int-eq2}齐次解）；方程\eqref{eq:int-eq2}有解当且仅当$\int_{0}^{1}g(t)x(t)\mathrm{d}t=0$（$x$为\eqref{eq:int-eq1}齐次解）。
\end{enumerate}
\end{examples}

\begin{definition}
设$\mathscr{X}$为$B$空间，对任意$T\in\mathscr{L}(\mathscr{X})$，定义
\begin{align*}
R(T)\triangleq T(\mathscr{X}),\quad N(T)\triangleq\{x\in\mathscr{X}\mid Tx=\theta\}.
\end{align*}
对任意$M\subseteq\mathscr{X}$，$N\subseteq\mathscr{X}^*$，定义
\begin{align*}
^\perp M\triangleq\{f\in\mathscr{X}^*\mid\langle f,x\rangle=0,\forall x\in M\},\quad N^\perp\triangleq\{x\in\mathscr{X}\mid\langle f,x\rangle=0,\forall f\in N\}.
\end{align*}
若$f\in\mathscr{X}^*$，$x\in\mathscr{X}$满足$\langle f,x\rangle=0$，简记为$f\perp x$.
\end{definition}

\begin{definition}[闭值域算子]
设$\mathscr{X}$为$B$空间，$T\in\mathscr{L}(\mathscr{X})$，若满足
\begin{align*}
R(T)=\overline{R(T)},
\end{align*}
则称$T$为\textbf{闭值域算子}。
\end{definition}

\begin{theorem}
设$\mathscr{X}$为$B$空间，$T\in\mathscr{L}(\mathscr{X})$，则$\sigma(T)=\sigma(T^*)$.
\end{theorem}
\begin{proof}

只需证$T^{-1}\in\mathscr{L}(\mathscr{X})\iff (T^*)^{-1}\in\mathscr{L}(\mathscr{X}^*)$.

{\heiti 必要性}: 由$(T^*)^{-1}=(T^{-1})^*$，结论显然成立。

{\heiti 充分性}: 设$(T^*)^{-1}\in\mathscr{L}(\mathscr{X}^*)$，则$(T^{**})^{-1}\in\mathscr{L}(\mathscr{X}^{**})$。
因$T=T^{**}\big|_{\mathscr{X}}$，故$T$为单射且$R(T)\subseteq\mathscr{X}$为闭集。
反证假设$R(T)\neq\mathscr{X}$，则存在$x_0\in\mathscr{X}\setminus R(T)$，$x_0\neq\theta$。由\hyperref[theorem:泛函分析--(哈恩-巴拿赫)Hahn-Banach定理]{Hahn-Banach定理}，存在$f\in\mathscr{X}^*$使得
\begin{align*}
f(x_0)=\|x_0\|,\quad f(x)=0,\ \forall x\in R(T).
\end{align*}
于是$0=f(Ty)=(T^*f)(y),\ \forall y\in\mathscr{X}$，即$T^*f=0$，得$f=\theta$，矛盾。
故$R(T)=\mathscr{X}$，即$T^{-1}\in\mathscr{L}(\mathscr{X})$.

\end{proof}

\begin{theorem}\label{theorem:紧算子与闭值域算子的关系}
设$\mathscr{X}$为$B$空间，$A\in\mathfrak{C}(\mathscr{X})$，$T=I-A$，则$T$为闭值域算子。
\end{theorem}
\begin{proof}

$N(T)$为$\mathscr{X}$的闭子空间，定义商空间上的算子
\begin{align*}
\widetilde{T}:\mathscr{X}/N(T)\to\mathscr{X},\quad \widetilde{T}[x]\triangleq Tx.
\end{align*}
则$R(\widetilde{T})=R(T)$，$\widetilde{T}$为有界线性单射，其逆算子存在。只需证$\widetilde{T}^{-1}$连续。

反证假设$\widetilde{T}^{-1}$不连续，则存在$[w_m]\to0$，$\widetilde{T}[w_m]\to0$，取子列满足$\|[w_{m_n}]\|\geqslant\varepsilon>0$。令$[x_n]=\frac{[w_{m_n}]}{\|[w_{m_n}]\|}$，则$\|[x_n]\|=1$，$\widetilde{T}[x_n]\to\theta$。
于是可取$x_n\in[x_n]$满足$\|x_n\|<2$，$(I-A)x_n\to\theta$。
由$A$为紧算子，存在子列$Ax_{n_k}\to z$，故$x_{n_k}=Ax_{n_k}+(I-A)x_{n_k}\to z$，得$Tz=\theta$，即$[z]=[\theta]$。
因此$\|[x_{n_k}]\|=\|[x_{n_k}-z]\|\leqslant\|x_{n_k}-z\|\to0$，与$\|[x_{n_k}]\|=1$矛盾。

\end{proof}

\begin{theorem}\label{theorem:Fredholm定理结论1}
设$\mathscr{X}$为$B$空间，$A\in\mathfrak{C}(\mathscr{X})$，$T=I-A$，若$N(T)=\{\theta\}$，则$R(T)=\mathscr{X}$.
\end{theorem}
\begin{proof}

反证假设$R(T)\neq\mathscr{X}$，定义$\mathscr{X}_0=\mathscr{X}$，$\mathscr{X}_k=T(\mathscr{X}_{k-1})\ (k\in\mathbb{N}^*)$，则
\begin{align*}
\mathscr{X}_0\supsetneqq\mathscr{X}_1\supsetneqq\mathscr{X}_2\supsetneqq\cdots.
\end{align*}
由\hyperref[lemma:Riesz引理]{Riesz引理}，存在$y_k\in\mathscr{X}_k$，$\|y_k\|=1$，满足$d(y_k,\mathscr{X}_{k+1})\geqslant\frac{1}{2}$。
对任意$p,n\in\mathbb{N}$，有
\begin{align*}
\|Ay_n-Ay_{n+p}\|=\|y_n-Ty_n+Ty_{n+p}-y_{n+p}\|\geqslant\frac{1}{2},
\end{align*}
其中$Ty_n-Ty_{n+p}+y_{n+p}\in\mathscr{X}_{n+1}$。
该式与$A$的紧性矛盾，故$R(T)=\mathscr{X}$.

\end{proof}

\begin{lemma}
设$\mathscr{X}$为$B$空间，$A\in\mathfrak{C}(\mathscr{X})$，$T=I-A$，则$\dim N(T)<\infty$，$\dim N(T^*)<\infty$.
\end{lemma}
\begin{proof}

令$\mathscr{X}_0=N(T)$，$B_1=\{x\in\mathscr{X}_0\mid\|x\|\leqslant1\}$，则$B_1=\overline{A(B_1)}$。
因$A$为紧算子，$B_1$为紧集，故$\dim\mathscr{X}_0=\dim N(T)<\infty$。
同理可证$\dim N(T^*)<\infty$.

\end{proof}

\begin{lemma}
设$\mathscr{X}$为$B$空间，$N(T)=\mathrm{span}\{x_1,x_2,\cdots,x_n\}$，则存在闭线性子空间$\mathscr{X}_1\subseteq\mathscr{X}$，使得
\begin{align*}
\mathscr{X}=\mathrm{span}\{x_1,x_2,\cdots,x_n\}\oplus\mathscr{X}_1.
\end{align*}
\end{lemma}
\begin{proof}

由\hyperref[theorem:泛函分析--(哈恩-巴拿赫)Hahn-Banach定理]{Hahn-Banach定理}，存在$g_1,g_2,\cdots,g_n\in\mathscr{X}^*$满足$g_i(x_j)=\delta_{ij}$，$1\leqslant i,j\leqslant n$。
令$\mathscr{X}_1=\bigcap\limits_{i=1}^n N(g_i)$，则$\mathscr{X}_1$为闭线性子空间，满足
\begin{enumerate}[(1)]
\item $\mathrm{span}\{x_1,x_2,\cdots,x_n\}\cap\mathscr{X}_1=\{\theta\}$;
\item 对任意$x\in\mathscr{X}$，取$c_i=g_i(x)$，有$x-\sum\limits_{i=1}^n c_ix_i\in\mathscr{X}_1$.
\end{enumerate}
因此$\mathscr{X}=\mathrm{span}\{x_1,x_2,\cdots,x_n\}\oplus\mathscr{X}_1$.

\end{proof}

\begin{lemma}
设$\mathscr{X}$为$B$空间，$N(T^*)=\mathrm{span}\{f_1,f_2,\cdots,f_m\}$，则存在$y_1,y_2,\cdots,y_m\in\mathscr{X}$，使得$f_i(y_j)=\delta_{ij}$，$1\leqslant i,j\leqslant m$.
\end{lemma}
\begin{proof}

定义线性连续映射$V:\mathscr{X}\to\mathbb{K}^m$，
\begin{align*}
V:x\mapsto(\langle f_1,x\rangle,\langle f_2,x\rangle,\cdots,\langle f_m,x\rangle).
\end{align*}
只需证$V$为满射。反证假设$V(\mathscr{X})$为$\mathbb{K}^m$的真子空间，由\hyperref[theorem:泛函分析--(哈恩-巴拿赫)Hahn-Banach定理]{Hahn-Banach定理}，存在$\alpha=(\alpha_1,\cdots,\alpha_m)\in\mathbb{K}^m\setminus\{0\}$使得$(\alpha,V(x))_{\mathbb{K}^m}=0,\ \forall x\in\mathscr{X}$，即
\begin{align*}
\left\langle\sum\limits_{j=1}^m\alpha_jf_j,x\right\rangle=0,\quad \forall x\in\mathscr{X}.
\end{align*}
得$\sum\limits_{j=1}^m\alpha_jf_j=\theta$，与$\{f_j\}$为$N(T^*)$的基矛盾。

\end{proof}

\begin{lemma}
设$\mathscr{X}$为$B$空间，$T\in\mathscr{L}(\mathscr{X})$，则$\overline{R(T)}=N(T^*)^\perp$.
\end{lemma}

\begin{proof}

(1) $\overline{R(T)}\subseteq N(T^*)^\perp$:
对任意$x\in\mathscr{X}$，$f\in N(T^*)$，有$f(Tx)=(T^*f)(x)=0$，故$R(T)\subseteq N(T^*)^\perp$。
因$N(T^*)^\perp$为闭集，故$\overline{R(T)}\subseteq N(T^*)^\perp$。

(2) $N(T^*)^\perp\subseteq\overline{R(T)}$:
反证假设存在$x_0\in N(T^*)^\perp\setminus\overline{R(T)}$，由Hahn-Banach定理，存在$f\in\mathscr{X}^*$满足
\begin{align*}
f(x_0)=\|x_0\|\neq0,\quad f(x)=0,\ \forall x\in\overline{R(T)}.
\end{align*}
于是$f(Tx)=0,\ \forall x\in\mathscr{X}$，即$T^*f=0$，故$f\in N(T^*)$。
结合$x_0\in N(T^*)^\perp$，得$f(x_0)=0$，矛盾。

\end{proof}

\begin{theorem}[Riesz-Fredholm定理]\label{theorem:泛函分析--Riesz-Fredholm定理}
设$\mathscr{X}$为$B$空间，$A\in\mathfrak{C}(\mathscr{X})$，$T=I-A$，则
\begin{enumerate}[(1)]
\item $\sigma(T)=\sigma(T^*)$;
\item $\dim N(T)=\dim N(T^*)<\infty$;
\item $R(T)=N(T^*)^\perp=\{x\in\mathscr{X}\mid f(x)=0,\forall f\in N(T^*)\}$，\quad $R(T^*)={}^\perp N(T)=\{f\in\mathscr{X}^*\mid f(x)=0,\forall x\in N(T)\}$.
\end{enumerate}
\end{theorem}
\begin{proof}

\textbf{情形1：$\dim N(T)=0$}
此时由前述定理得$R(T)=\mathscr{X}$，故$T$为双射。由Banach逆算子定理，$T^{-1}\in\mathscr{L}(\mathscr{X})$，即$0\notin\sigma(T)$。
又$\sigma(T)=\sigma(T^*)$，故$T^*$也为双射，即$\dim N(T^*)=0=\dim N(T)$，且
\begin{align*}
R(T^*)=\mathscr{X}^*={}^\perp N(T),\quad R(T)=\mathscr{X}=\{\theta\}^\perp=N(T^*)^\perp.
\end{align*}

\textbf{情形2：$\dim N(T)=n>0$}
设$N(T)=\mathrm{span}\{x_1,\cdots,x_n\}$，$N(T^*)=\mathrm{span}\{f_1,\cdots,f_m\}$。
先证$n=m$：
假设$n<m$，结合空间直和分解，构造算子$\widehat{T}$为满射，可推得$y_m\notin R(\widehat{T})$，矛盾，故$\dim N(T^*)\leqslant\dim N(T)$。
同理$\dim N(T^{**})\leqslant\dim N(T^*)$，又$\dim N(T)\leqslant\dim N(T^{**})$，故$\dim N(T)=\dim N(T^*)$。

结合闭值域算子定理与$\overline{R(T)}=N(T^*)^\perp$，得
\begin{align*}
R(T)=\overline{R(T)}=N(T^*)^\perp=\{x\in\mathscr{X}\mid f(x)=0,\forall f\in N(T^*)\}.
\end{align*}
由$\dim N(T^{**})=\dim N(T^*)=\dim N(T)$及$N(T)\subseteq N(T^{**})$，得$N(T^{**})=N(T)$，故
\begin{align*}
R(T^*)=\overline{R(T^*)}=N(T^{**})^\perp={}^\perp N(T).
\end{align*}
结论(1)由前述有界算子谱性质直接成立。

\end{proof}

\begin{theorem}
设$\mathscr{X}$为$B$空间，$A\in\mathfrak{C}(\mathscr{X})$，$T=I-A$，则存在闭线性子空间$\mathscr{X}_1$、有限维子空间$\mathscr{Y}_1$满足$\dim\mathscr{Y}_1=\dim N(T)$，使得
\begin{align*}
\mathscr{X}=N(T)\oplus\mathscr{X}_1=\mathscr{Y}_1\oplus R(T).
\end{align*}
\end{theorem}

\begin{definition}
设$\mathscr{X}$为$B$空间，$M\subseteq\mathscr{X}$为闭线性子空间，定义$M$的\textbf{余维数}为
\begin{align*}
\mathrm{codim}\,M\triangleq\dim(\mathscr{X}/M).
\end{align*}
\end{definition}
\begin{remark}
由\refthe{theorem:闭线性子空间的商空间完备}知$\mathscr{X}/M$也是$B$空间.
\end{remark}

\begin{theorem}
设$\mathscr{X}$为$B$空间，$A\in\mathfrak{C}(\mathscr{X})$，$T=I-A$，则
\begin{align*}
\dim N(T)=\mathrm{codim}\,R(T)<\infty.
\end{align*}
\end{theorem}
\begin{proof}

由\hyperref[theorem:泛函分析--Riesz-Fredholm定理]{Riesz-Fredholm定理}立得.
\end{proof}

\begin{example}
设$\mathscr{X}$是$B$空间，$M \subseteq \mathscr{X}$是一个闭线性子空间，$\mathrm{codim}M = n$，求证：存在线性无关集$\{\varphi_k\}_{k=1}^{n} \subseteq \mathscr{X}^{*}$，使得
\begin{align*}
M = \bigcap\limits_{k=1}^{n} N(\varphi_k).
\end{align*}
\end{example}

\begin{example}
设$\mathscr{X},\mathscr{Y}$是两个$B$空间，$T \in \mathscr{L}(\mathscr{X},\mathscr{Y})$是满射的。定义$\widetilde{T}: \mathscr{X}/N(T) \to \mathscr{Y}$如下：
\begin{align*}
\widetilde{T}[x] = Tx \quad (\forall x \in [x])(\forall [x] \in \mathscr{X}/N(T)).
\end{align*}
求证：$\widetilde{T}$是线性同胚映射。
\end{example}

\begin{example}
设$\mathscr{X}$是$B$空间，$M,N_1,N_2$都是$\mathscr{X}$的闭线性子空间，如果
\begin{align*}
M \oplus N_1 = \mathscr{X} = M \oplus N_2,
\end{align*}
求证：$N_1$和$N_2$同胚。

提示 只要证明$N_1,N_2$都与$\mathscr{X}/M$同胚。
\end{example}

\begin{example}
设$A \in \mathfrak{C}(\mathscr{X})$，$T = I - A$，求证：
\begin{enumerate}[(1)]
\item $\forall [x] \in \mathscr{X}/N(T), \exists x_0 \in [x]$，使得$\|x_0\| = \|[x]\|$；
\item 若$y \in \mathscr{X}$，使方程$Tx = y$有解，则其中必有一个解达到范数最小。
\end{enumerate}
\end{example}

\begin{example}
设$A \in \mathfrak{C}(\mathscr{X})$，且$T = I - A, \forall k \in \mathbb{N}$，求证：
\begin{enumerate}[(1)]
\item $N(T^k)$是有穷维的；
\item $R(T^k)$是闭的。
\end{enumerate}
\end{example}

\begin{example}
设$M$是$B$空间$\mathscr{X}$的闭线性子空间，称满足$P^2 = P$（幂等性）的由$\mathscr{X}$到$M$上的一个有界线性算子$P$为由$\mathscr{X}$到$M$上的投影算子。求证：
\begin{enumerate}[(1)]
\item 若$M$是$\mathscr{X}$的有穷维线性子空间，则必存在由$\mathscr{X}$到$M$上的投影算子；
\item 若$P$是由$\mathscr{X}$到$M$上的投影算子，则$I-P$是由$\mathscr{X}$到$R(I-P)$上的投影算子；
\item 若$P$是由$\mathscr{X}$到$M$上的投影算子，则$\mathscr{X} = M \oplus N$，其中$N = R(I-P)$；
\item 若$A \in \mathfrak{C}(\mathscr{X})$，且$T = I-A$，则在代数与拓扑同构意义下，$N(T) \oplus \mathscr{X}/N(T) = \mathscr{X} = R(T) \oplus \mathscr{X}/R(T)$。
\end{enumerate}
\end{example}








\end{document}